\documentclass{practice}

\title{6}
\date{\today}

%\usepackage{etoolbox}
%\makeatletter
%\preto{\@verbatim}{\topsep=0pt \partopsep=-2\parsep }
%\makeatother

\usepackage{crysymb}

\usepackage{fancyvrb}
\fvset{listparameters=\setlength{\topsep}{0pt}\setlength{\partopsep}{0pt}}

\usepackage{listings}

\lstset{
basicstyle=\small\ttfamily,
columns=flexible,
breaklines=true
}

\begin{document}
\maketitle

\begin{task}{OpenSSL}
  \textit{Preface.}
  For working with public-key cryptography keypairs, OpenSSL provides the following tools:
  \begin{itemize}
    \item \href{https://docs.openssl.org/master/man1/openssl-genpkey/}{\texttt{genpkey}}: utility for generating a private key (RSA, EC) or parameters (DH, DSA, EC)
    \begin{itemize}
      \item Alternatively: \href{https://docs.openssl.org/master/man1/openssl-genrsa/}{\texttt{genrsa}}, \href{https://docs.openssl.org/master/man1/openssl-dhparam/}{\texttt{dhparam}}, \href{https://docs.openssl.org/master/man1/openssl-ecparam/}{\texttt{ecparam}}
    \end{itemize}
    \item \href{https://docs.openssl.org/master/man1/openssl-pkey/}{\texttt{pkey}}: utility for processing keys, e.g. change their format (PEM, DER) or passphrase
    \begin{itemize}
      \item Alternatively: \href{https://docs.openssl.org/master/man1/openssl-rsa/}{\texttt{rsa}}, \href{https://docs.openssl.org/master/man1/openssl-ec}{\texttt{ec}}
    \end{itemize}
    \item \href{https://docs.openssl.org/master/man1/openssl-pkeyutl/}{\texttt{pkeyutl}}: utility for actual operations with keys (e.g. encryption or signing)
    \begin{itemize}
      \item \href{https://docs.openssl.org/master/man1/openssl-rsautl/}{\texttt{rsautl}} is deprecated, do not use it
    \end{itemize}
  \end{itemize}

  To generate parameters instead of a private key with the \texttt{genpkey} utility, you must use the \texttt{-genparam} option.
  For example, if you one day feel the strong urge to generate a large prime before going to bed, you can use the following command:
  \begin{Verbatim}
openssl genpkey -genparam -algorithm dh \
    -pkeyopt dh_paramgen_prime_len:3072
  \end{Verbatim}

  If you wonder what the dots, plus signs, and asterisks that appear mean, then for DH there is a nice InfoSec Stack Exchange \href{https://security.stackexchange.com/a/140639}{\textit{post}} on it.

  In practice, perhaps the primary use case for generating parameters in this manner with OpenSSL is to publish the parameters or share them with other entities.
  These parameters then form the basis of subsequent procedures\footnotemark{}.
  For example, to generate a private key with OpenSSL based on a set of parameters, you can use the \texttt{-paramfile} option.
  \footnotetext{E.g. the frowned upon approach of generating custom Diffie-Hellman parameters for TLS 1.2.}

  \textit{Task.}
  For this task, familiarise yourself with the OpenSSL documentation.
  Then, generate a password-protected P-384 key encrypted with AES-128 in CBC mode.\footnotemark{}
  \footnotetext{\texttt{genpkey} does not support AEAD for some reason, and AES in CTR mode has an implementation bug for this use.
  Marvels of OpenSSL\dots{}.}
\end{task}

\begin{task}{HPKE}
  \textit{Preface.}
  In the lecture we briefly discussed IES and ECIES.
  However, as you may recall, the various (!) ECIES specifications are not freely accessible, and IES itself does not have a single \enquote{canonical} standard.
  The modern alternative is \emph{Hybrid Public Key Encryption} (HPKE), specified in \href{https://datatracker.ietf.org/doc/html/rfc9180}{RFC 9180}.

  In its commonly used form, HPKE combines a KEM, a KDF, and an AEAD scheme.
  To encrypt a message, the sender must know the recipient's public key -- whether long-term or ephemeral (e.g. to achieve forward secrecy).

  You may then ask, should the sender be authenticated as well?
  HPKE supports four modes of operation:
  \begin{itemize}
    \item Base -- The sender generates a fresh ephemeral keypair for each message and remains anonymous (at the scheme level) to the recipient.
    \item PSK -- \emph{Pre-shared key}.
    A symmetric key known to both parties is incorporated into the key \enquote{schedule}.
    This provides symmetric authenticity guarantees.
    \item Auth -- The sender additionally uses a static (long-term) key pair to authenticate themselves to the recipient (we will study public-key authenticity later).
    \item Auth + PSK -- Auth mode with an additional PSK mixed into the key schedule.
  \end{itemize}

  \textit{Task.}
  pyca/cryptography has \href{https://github.com/pyca/cryptography/pull/14126}{\emph{recently}} implemented HPKE (in base mode).
  in fact, it will not be available before version 47 (the current latest version is 45.0.5).

  However, the \href{https://pycryptodome.readthedocs.io/en/latest/src/protocol/hpke.html}{\emph{pycryptodome}} library already supports HPKE.

  For this task do the following:
  \begin{enumerate}
    \item Acquire the server's (Discord bot's) long-term P-384 EC public key.
    \item Generate your own P-384 ephemeral keypair.
    \item Submit your ephemeral public key to the bot, and receive a ciphertext.
    \item Decrypt the AES-128-ECB encrypted ciphertext with the context (info)
    \begin{Verbatim}
ITC8280 week 6
    \end{Verbatim}
    The server is using HPKE in base mode.
  \end{enumerate}
\end{task}
\end{document}
